\documentclass[10pt,letterpaper]{article}
\usepackage[margin=0.62in]{geometry}
\usepackage[T1]{fontenc}
\usepackage[default]{sourcesanspro}
\usepackage{enumitem}
\usepackage{titlesec}
\usepackage[colorlinks=true,urlcolor=mutedblue]{hyperref}
\usepackage{xcolor}
\definecolor{mutedblue}{HTML}{51758B}
\usepackage{tabularx}
\usepackage{microtype}
\pagestyle{empty}
\setlength{\parindent}{0pt}
\setlength{\parskip}{0pt}
\setlist[itemize]{leftmargin=1.15em,topsep=3pt,itemsep=2pt,parsep=0pt,partopsep=0pt}
\titleformat{\section}{\bfseries\normalsize\uppercase}{}{0pt}{}[\titlerule]
\titlespacing*{\section}{0pt}{13pt}{6pt}
\newcommand{\entry}[2]{\noindent\begin{tabularx}{\textwidth}{@{}Xr@{}}\textbf{#1} & #2\\\end{tabularx}}
\begin{document}
\begin{center}
{\Large\bfseries Ziyi Weng}\\[4pt]
\href{mailto:ziyiweng39@gmail.com}{ziyiweng39@gmail.com}\quad|\quad
\href{https://github.com/ziyiweng}{GitHub}\quad|\quad
\href{https://www.linkedin.com/in/ziyiweng}{LinkedIn}
\end{center}
\vspace{-6pt}

\section{Education}
\entry{University of Illinois Urbana-Champaign}{Expected Graduation: May 2030}
Bachelor of Science in Mathematics

\section{Technical Skills}
\textbf{Languages:} Python, C++, Rust, Java\\
\textbf{Frameworks \& Libraries:} PyTorch, NumPy, pandas, Matplotlib\\
\textbf{Developer Tools:} Git, LaTeX

\section{Relevant Coursework}
\textbf{Fall 2026:} CS 128 (Introduction to Computer Science II), CS 128H (Honors), MATH 241 (Calculus III), MATH 314 (Introduction to Higher Mathematics)\\
\textbf{Spring 2027 (Planned):} CS 225 (Data Structures), CS 233 (Computer Architecture), MATH 416 (Abstract Linear Algebra), STAT 400 (Statistics and Probability I)

\section{Research \& Projects}
\entry{\href{https://github.com/uiuc-2026-code-ada-team10/will-it-fit}{Will-it-fit} $|$ Code Ada 2026 Hackathon}{Oct 2026}
\begin{itemize}
\item Built an interactive room design tool using Cursor for AI-assisted development, allowing users to turn 2D floor plans into customizable 3D room layouts.
\item Defined feature requirements and iteratively tested and refined 3D furniture customization, object manipulation, and boundary collision alerts.
\end{itemize}
\vspace{7pt}

\entry{ConvLSTM-Based ENSO Forecasting}{Mar 2026 -- Jul 2026}
\textit{Co-First Author $|$ Paper accepted at OCEANS 2026 (IEEE)}
\begin{itemize}
\item Designed a 49K-parameter ConvLSTM pipeline in PyTorch with wavelet bandpass filtering, spatial self-attention, and PDE-based regularization for long-range ENSO forecasting.
\item Matched the 18-month forecasting performance of Ham et al. (Nature, 2019) with 40$\times$ fewer parameters and no climate-simulation pretraining.
\end{itemize}
\vspace{7pt}

\entry{Mathematics and Statistics Research Challenge}{Jun 2025 -- Aug 2025}
\textit{Individual Participant $|$ Distinction Award (Top 4\%)}
\begin{itemize}
\item Independently investigated a combinatorial optimization problem involving minimum covering sets on a 3D lattice.
\end{itemize}

\section{Honors \& Awards}
\entry{Silver Medal, New Zealand Maths Olympiad}{2024, 2025}
\par\vspace{8pt}\noindent
\entry{NZMOC Mathematical Olympiad Training Camp}{2024}
\begin{itemize}
\item Selected among approximately 25--30 students nationwide for the IMO team selection process.
\end{itemize}

\section{Leadership \& Activities}
\entry{Jane Street WiSE Program}{Feb 2026}
\textit{Participant}
\begin{itemize}
\item Explored applications of probability, statistics, and computer science in quantitative trading.
\end{itemize}
\vspace{7pt}

\entry{Inter-School Mathematics Olympiad}{Feb 2025 -- Apr 2025}
\textit{Founder \& Organizer}
\begin{itemize}
\item Organized my high school's first participation in ISMO and evaluated 150+ competition submissions from students across New Zealand.
\end{itemize}
\end{document}
